\section{One amplification block: words and overlap}\label{sec:block}

We now construct one longer pair from the preceding oracle pair.  This section defines the nested words and proves the consecutive-power overlap needed for their side runs.  Section~\ref{sec:layout} constructs the observation of those words, and Section~\ref{sec:coefficients} proves its score estimate.

\begin{proposition}[One amplification block]\label{prop:block}
For the parameters of Section~\ref{sec:scales}, let $1\le i\le D$ and suppose a level-$(i-1)$ pair satisfies Definition~\ref{def:invariant} with constants $C,\omega,M_B$.  Set
\[
 C'=\max\{C,C_{\rm abs}/\omega\}.
\]
The input pair remains valid when $C$ is replaced by $C'$, and there exists a level-$i$ pair satisfying the definition with the same constants $C',\omega,M_B$, with
\[
 d_i=d_{i-1}+(b_i-1)/4-.05,
 \qquad E_i\le\Delta+(\ell_i/\ell_{i-1})E_{i-1}.
\]
Every preceding level constructed with $C$ also remains valid with $C'$, and this enlarged value is retained in later blocks.  The conclusion is uniform over all levels along the stated joint asymptotic regime.
\end{proposition}

Fix $i$ for the proof.  The value $C'$ meets the hypothesis of Lemma~\ref{lem:lift}.  Every preceding upper count bound remains true after replacing $C$ by $C'$, since its exponent $2Bh+M_B$ at a level $h\ge0$ is nonnegative; the lower count bounds, scores, and projection errors do not involve $C$.  The enlarged constant still depends only on $\delta$ and satisfies $\log C'=O(\Lambda)$.  We henceforth write $C$ for $C'$ and use this common enlarged constant in the counts constructed below and in later blocks.

Abbreviate
\[
\begin{gathered}
 b=b_i,\quad N=N_i,\quad d=d_{i-1},\quad d'=d_i,\\
 x=x_{i-1},\quad y=y_{i-1},\quad \ell=\ell_{i-1},\quad k=k_{i-1}.
\end{gathered}
\]
Within this block write $a_0=a_0(b)$ for the contraction exponent defined in \eqref{eq:contraction-exponent}.  Thus $\ell\in[N^{a_0}/2,N^{a_0}]$.  Since $r_i\le r_{i-1}$, the old score satisfies
\begin{equation}\label{eq:old-score}
 \left\|\frac{d\nu_{i-1}}{d\mu_{i-1}}-1\right\|_{L^{r_i}(\mu_{i-1})}^2
       \le N^{-a_0d}.
\end{equation}

\subsection{Nested signed words}

The two new words will use the same nested motif at $b$ depths.  Choose the fixed degree cutoff and the cancellation order
\[
 r_*=2^{10B+40},\qquad
 s_*=\left\lceil100(d_D+r_*+B+10)/\epsilon\right\rceil,
 \qquad L=2^{s_*}.
\]
Let $\mathcal I=\{0,\ldots,L-1\}$ and define the signs by
\[
 \sum_{a\in\mathcal I}\sigma_a z^a
       =\prod_{v=0}^{s_*-1}(1-z^{2^v}).
\]
Extend $\sigma_a$ by zero outside $\mathcal I$, and put $\chi={\bf1}_{\mathcal I}$.  The order $s_*=O_B((D+1)^2)$ is large enough for the signed smoothing below at every block.  Lemma~\ref{lem:scale-margins} gives, uniformly,
\begin{equation}\label{eq:block-losses}
\begin{gathered}
 L^{O_B(r_*B)}=N^{o(\epsilon)},\\
 \exp\bigl(O_B((s_*+1)^2+(s_*+1)(D+1)\Lambda)\bigr)
       =N^{o(\epsilon)}.
\end{gathered}
\end{equation}
Here and below fixed-depth constants may depend on $B$.

Start with $U_0^+=x$, $U_0^-=y$, both of length $S_0=\ell$, and with overlap count $h_0=k$.  Once a positive integer $P_j$ has been chosen, form $U_j^\tau$, for $\tau\in\{+1,-1\}$, from $P_j$ consecutive level-$(j-1)$ atoms.  Use type $+$ except at the labels
$\lfloor P_j/2\rfloor+a$, $a\in\mathcal I$, where the type is $\tau\sigma_a$.  We choose the counts below so that these labels lie in $[1,P_j]$.  The common length is
\[
 S_j=P_jS_{j-1}.
\]
The two words are distinct: at a motif label their atom types differ, and equal-length atoms give a fixed parse back to the distinct old words.

\subsection{Word sizes, windows, and paired counts}

Write $N^{a_j}$ for the target length of a level-$j$ word.  It then contains about $N^{a_j-a_{j-1}}$ lower atoms.  An overlap chunk will contain $h_{j-1}$ such atoms, with $h_{j-1}$ within the controlled channel factor of their word length $N^{a_{j-1}}$.  To leave about $N^\epsilon$ paired chunks, choose the remaining exponents by
\[
 a_j(b)=2a_{j-1}(b)+\epsilon\quad(1\le j\le b).
\]
This recurrence and \eqref{eq:contraction-exponent} give the identities
\begin{equation}\label{eq:local-exponents}
 a_j(b)=(1+\epsilon)2^{j-b}-\epsilon\quad(0\le j\le b),
 \qquad a_b(b)=1.
\end{equation}
In the rest of this block, $a_j$ abbreviates $a_j(b)$.

A level-$j$ child window will contain $T_j$ level-$(j-1)$ atoms.  Define
\[
 \kappa_j=a_{j-1}-3(2^j-1)\epsilon,
 \qquad T_j=4\left\lfloor N^{\kappa_j}/4\right\rfloor
       \quad(1\le j\le b).
\]
The $\kappa_j$ are positive for all sufficiently large instances, uniformly in the block.  Directly from \eqref{eq:local-exponents},
\begin{equation}\label{eq:kj}
\begin{aligned}
 2\kappa_j-\kappa_{j+1}-\epsilon&=\epsilon\quad(1\le j<b),\\
 \kappa_b&>\epsilon,\qquad
 \kappa_j<a_{j-1}<a_j-a_{j-1}\quad(1\le j\le b).
\end{aligned}
\end{equation}
The first identity makes $T_j$ larger than the typical displacement from $T_{j+1}$ collections of about $N^\epsilon$ fair bits: the squared window width divided by that bit count has order $N^\epsilon$.  The remaining inequalities leave room for the windows inside their words.  We also have
\begin{equation}\label{eq:power_gap}
 \sum_{j=1}^b\kappa_j<1,
\end{equation}
since $\kappa_j<a_{j-1}<2^{j-1-b}$; this will bound the number of vertices in the observation tree.

For $1\le j<b$, choose
\[
\begin{aligned}
 P_j&=\left\lfloor N^{a_j}/S_{j-1}\right\rfloor,
 &h_j&=\lceil C h_{j-1}P_j\rceil,\\
 t_j&=2\left\lfloor P_j/[100(2h_{j-1}+1)]\right\rfloor.
\end{aligned}
\]
For the root, choose
\[
\begin{aligned}
 J&=\left\lfloor\frac{\lfloor N/S_{b-1}\rfloor-T_b}
                         {2h_{b-1}+1}\right\rfloor,
 &P_b&=J(2h_{b-1}+1)+T_b,\\
 h_b&=\lceil C h_{b-1}P_b\rceil.
\end{aligned}
\]
The counts $t_j$ and $J$ are the numbers of paired chunks used for local and outer displacements, respectively.  The factors $2$ and $4$ in their definitions and in $T_j$ ensure the parity and divisibility needed by the partition.

We record the size consequences before using the words.  For $j<b$, flooring loses at most $S_{j-1}=o(N^{a_j})$, so
\[
 S_j\in[N^{a_j}/2,N^{a_j}],\qquad
 P_j\asymp N^{a_j-a_{j-1}},\qquad T_j=o(P_j).
\]
The overlap counts remain within the declared uniform bound.  Indeed $h_j\ge S_j$ by induction, and
\[
 \frac{h_j}{S_j}\le C\frac{h_{j-1}}{S_{j-1}}+1
       \le C^2\frac{h_{j-1}}{S_{j-1}}
\]
because $h_{j-1}/S_{j-1}\ge1$ and $C>2$.  Starting from the level-$(i-1)$ bound gives
\begin{equation}\label{eq:nested-counts}
 S_j\le h_j\le C^{2B(i-1)+M_B+2j}S_j
             \le C^{2Bi+M_B}S_j\quad(0\le j\le b),
\end{equation}
In particular the extra count factor is at most $\exp(O_B((D+1)\Lambda))$.

Using these bounds in the definitions of $t_j$ and $J$ yields
\begin{equation}\label{eq:scale_t}
 e^{-O_B((D+1)\Lambda)}N^\epsilon\le t_j\le O(N^\epsilon)
       \quad(1\le j<b),
\end{equation}
and the same lower and upper bounds for $J$.  All these counts tend to infinity uniformly.  At the root, $J=o(T_b)$ by $\kappa_b>\epsilon$, and
\[
 N-S_b=O((h_{b-1}+1)S_{b-1})
 \le N^{2a_{b-1}+o(\epsilon)}
 =N^{1-\epsilon+o(\epsilon)}=o(N).
\]
Here the count overhead is $N^{o(\epsilon)}$ by \eqref{eq:uniform-absorption}, and $2a_{b-1}=1-\epsilon$ by \eqref{eq:local-exponents}.  Thus also $S_b\in[N/2,N]$.  Working backward from $\kappa_b>\epsilon$ in \eqref{eq:kj} gives $\kappa_j>\epsilon$ at every depth.  The uniform loss bound \eqref{eq:block-losses} therefore gives $L=o(T_j)$ and $L=o(P_j)$ uniformly, so all motif labels used above fit.  We take
\[
 x_i=U_b^+,\qquad y_i=U_b^-,\qquad \ell_i=S_b,\qquad k_i=h_b.
\]

The following estimate quantifies the smoothing supplied by a hidden binomial offset.  It will be used in the early overlap comparison and in the main score proof.  In those applications, its local parameter $m$ is one of the paired counts $J$ or $t_j$, while $g$ is the actual number of hidden orientation trials.
\begin{lemma}[Signed binomial smoothing]\label{lem:binomial-smoothing}
Let $\sigma$ be the Prouhet--Thue--Morse sequence of order $s_*$ above.  Suppose $\pi$ is an integer translation or reflection of a fair binomial probability mass function with $g\ge c_0\omega^2m$ trials, where $c_0>0$ is absolute and $m\to\infty$ as in the construction.  Then
\begin{align}
 \|\sigma*\pi\|_\infty
 &\le \exp(O((s_*+1)^2+(s_*+1)\Lambda))\,
       m^{-(s_*+1)/2},\nonumber\\
 \|\sigma*\pi\|_2^2
 &\le \exp(O((s_*+1)^2+(s_*+1)\Lambda))\,
       m^{-s_*-1/2},\qquad
 \|\pi\|_\infty\le O(\omega^{-1}m^{-1/2}).\label{eq:sig_conv}
\end{align}
The norms use counting measure on the integers, and the implicit constants may depend on $c_0$.
\end{lemma}

\begin{proof}
Translations and reflections do not affect the estimates.  For the first two, use Fourier inversion and Parseval.  On $[-\pi,\pi]$ the sign transform has modulus at most $2^{s_*(s_*-1)/2}|z|^{s_*}$, and the binomial transform has modulus at most $e^{-cgz^2}$.  Gaussian integration contributes, besides the power of $g$, a factor $\exp(O((s_*+1)\log(s_*+1)))$.  The last estimate is the binomial maximum bound from Stirling's formula.
\end{proof}

\subsection{Overlap for each nested atom}\label{sec:side-overlap}

\begin{lemma}[Overlap of the nested atoms]\label{lem:nested-overlap}
Under the input hypotheses of Proposition~\ref{prop:block}, the words and counts constructed above satisfy, for every $1\le j\le b$,
\[
 \ov\bigl(\Tr((U_j^+)^{h_j}),\Tr((U_j^+)^{h_j+1})\bigr)\ge\omega.
\]
\end{lemma}

\begin{proof}
The old score and the two projection errors give
\begin{equation}\label{eq:old_H}
 H^2(\Tr(x),\Tr(y))\le O(N^{-a_0d}+E_{i-1}).
\end{equation}
Indeed the score bounds the squared Hellinger distance of the old oracle laws, and contraction followed by the Hellinger triangle inequality handles the two projection errors.

For the induction on $j$, assume the overlap statement at level $j-1$; at $j=1$ it is the input invariant.  Lemma~\ref{lem:lift}, with base word $U_{j-1}^+$, overlap count $h_{j-1}$, replication $P_j$, and $F=h_j\ge C h_{j-1}P_j$, gives overlap at least $2/3$ for the two pure comparison words formed by repeating $P_j$ copies of $U_{j-1}^+$.  Replacing one such length-$S_j$ unit by $U_j^+$ changes at most $L$ lower atoms.  For $h_j$ or $h_j+1$ repeated units, the squared Hellinger replacement cost is therefore at most
\begin{equation}\label{eq:nested-replacement}
 (h_j+1)L\,H^2(\Tr(U_{j-1}^+),\Tr(U_{j-1}^-)).
\end{equation}

The two phases of the depth schedule give the margins needed here.  From \eqref{eq:local-exponents} and the choice of $\epsilon$, for $b=2$ and $d\ge2.20$,
\begin{equation}\label{eq:b2gap}
 a_0d-a_1>.03,\qquad 2a_0d>1.06.
\end{equation}
For $i>j_0$, we have $b=B$ and $d\ge2^B+2.20$, giving
\begin{equation}\label{eq:bbgap}
 a_0d>1+2^{-B}.
\end{equation}
For example, before the $\epsilon d$ correction the first two left sides are at least $.05$ and $1.10$, and the last is at least $1+2.20\,2^{-B}$.  The bound $d\le2.20+(B/4)D$ and $\epsilon=2^{-4B-100}/(D+1)$ leave the displayed margins uniformly.

In a block with $i>j_0$, the two level-$(j-1)$ atoms differ on at most $L^{j-1}$ old atoms.  Tensorization, concatenation, and \eqref{eq:old_H} therefore bound \eqref{eq:nested-replacement} by $o(1)$: use \eqref{eq:bbgap}, $h_j\le e^{O_B((D+1)\Lambda)}N$, $L^b=N^{o(\epsilon)}$, and $E_{i-1}\le n(D+1)\Delta$.  In a preliminary block, $b=2$.  At $j=1$ the same argument uses $h_1\le e^{O_B((D+1)\Lambda)}N^{a_1}$ and the first inequality in \eqref{eq:b2gap}.  At $j=2$ we need the sharper comparison
\begin{equation}\label{eq:side_refine}
 H^2(\Tr(U_1^+),\Tr(U_1^-))
 \le N^{-2a_0d+o(\epsilon)}+O(P_1E_{i-1}+\Delta),\qquad b=2.
\end{equation}
The second inequality in \eqref{eq:b2gap} then makes \eqref{eq:nested-replacement} $o(1)$ at the root as well.

To prove \eqref{eq:side_refine}, report a single window of $T_1$ separate old oracle outputs.  Its nominal start is after
$A=\lfloor(P_1-T_1)/2\rfloor$ old atoms.  Put $t_1$ paired pure-$x$ chunks before and after the window, with counts $h_0+E,h_0+1-E$ for independent fair bits $E$.  Use the paired flag construction at the input overlap.  Complete the two sides by pure-$x$ residuals of counts
\[
 A-t_1(h_0+1/2),\qquad P_1-T_1-A-t_1(h_0+1/2).
\]
These are nonnegative integers.  Chunks and residuals have fixed relative IDs and each reports one unseparated ordinary trace, together with its flag and any revealed bit.  The window start displacement $V$ is the sum of the $t_1$ bits minus $t_1/2$, so $|V|\le t_1/2=o(T_1)$.  Every motif label therefore lies in the window.

Conditional on the orientation bits, sample the $T_1$ window outputs independently of one another and of the entire collection of side flags and emissions, using $\mu_{i-1}$ for an $x$ atom and $\nu_{i-1}$ for a $y$ atom at the corresponding rank.  Inside the window, report the old oracle outputs with their relative ranks, but not their absolute start or type labels.  Apply $\mathcal K_{i-1}$ independently to these outputs and concatenate in the scheduled order.  Given the orientation bits, averaging the flags produces the exact independent trace laws of the other pieces.  The common projection therefore approximates the genuine trace with error at most $T_1E_{i-1}$ under either hypothesis.

Except on side metadata of probability $o(\Delta)$, at least $\omega^2t_1/2$ orientations are hidden, by binomial concentration and \eqref{eq:scale_t}.  Conditional on such metadata, $V$ is a translated fair binomial with that many trials.  Relative to $Q=\mu_{i-1}^{\otimes T_1}$, let $G_+,G_-$ be the conditional densities of the window outputs.  For fixed $V$, each is a product of at most $L$ factors $1+Z$ at distinct positions, averaged over $V$, where $Z$ is the old score.  Lemma~\ref{lem:moment} gives for its homogeneous degree-$r$ part the squared $L^{r_i}(Q)$ bound
\[
 (C_1r_i)^{2r}N^{-a_0dr}\binom Lr.
\]
Convexity permits averaging over $V$ at the same bound.  Geometric summation gives
\[
 \|G_\pm-1\|_{r_i}\le N^{-a_0d/2+o(\epsilon)},
\]
and the squared norm of the terms of degree at least two is at most $N^{-2a_0d+o(\epsilon)}$.  In $G_+-G_-$ the singleton coefficient array is, up to signs, the convolution of $\sigma$ with the law of $V$.  Equation~\eqref{eq:sig_conv} and \eqref{eq:scale_t} bound its square sum strongly enough that
\[
 \|G_+-G_-\|_{r_i}^2\le N^{-2a_0d+o(\epsilon)};
\]
indeed $t_1^{-s_*}\le N^{-\epsilon s_*}
\exp(O_B((D+1)s_*\Lambda))$.

On typical metadata, the region $G_+<1/2$ has mass $o(\Delta)$ under both laws by the same moment argument as in Lemma~\ref{lem:clipping}.  On its complement, the squared density difference bounds Hellinger distance.  Adding the atypical metadata and projection errors proves \eqref{eq:side_refine}, with fixed numerical factors absorbed in $o(\epsilon)$.  Thus all replacement costs in \eqref{eq:nested-replacement} are $o(1)$.  The pure overlap $2/3$ loses at most the two replacement total variations, and remains above $\omega<1/8$.  This completes the overlap induction.
\end{proof}

The words, counts, and overlaps required for the observation are now available at every nested depth.  The next section uses their paired side runs to hide the relative locations of the motifs.
